\chapter{EhBASIC}\label{cha:basic}

The machine speaks EhBASIC~2.22 --- Lee Davison's Enhanced BASIC --- with
this machine's additions: graphics statements that ride the blitter,
floating point on the MATH unit, and an escape hatch to the shell. It has
no sound \emph{keyword}: a program that wants a note poke\-s the sound
sequencer at \texttt{\$D5E0} or the OPL2 at \texttt{\$D480} directly,
which on a machine with a friendly memory map is not much of a hardship
(Chapter~\ref{cha:io}).

\begin{type}
RUN EHBASIC
\end{type}

It comes up ready: the banner, \texttt{47103 Bytes free} --- more than
twenty-one C64s more than a C64 --- and a
\verb|Ready| prompt --- the traditional \texttt{Memory size ?} question is
gone; the machine knows.

\section{Ten minutes of it}
\begin{type}
RUN "DEMOS.BAS"
\end{type}

\screen{shots/demos.png}{The demo menu. A key picks an entry; each demo
returns here when it ends.}

The menu chains: \verb|RUN "name"| (or \verb|LOAD| inside a running
program) loads another BASIC program and runs it --- that is how one BASIC
program hands over to another, exactly as on a C64.

\section{The machine's own statements}
\begin{type}
GRAPHICS 2          640x480 bitmap over the text
PLOT X,Y,C          LINE X1,Y1,X2,Y2,C
TRI X1,Y1,X2,Y2,X3,Y3,C
PALETTE I,R,G,B     GCLS            GRAPHICS 0
\end{type}

Colours 0--15 belong to the text screen; demos use 16 and up. \verb|GRAPHICS 0|
puts the text mode and its palette back, whatever the program changed.

\section{Hardware sprites}
The 128 sprites VICKY carries are reachable from BASIC, which is unusual
enough to be worth saying plainly: this is not a bitmap being redrawn, it
is the video chip carrying the picture for you.

\begin{type}
SPRDEF n,page,w,h,bpp   give sprite n a shape
SPRITE n,x,y            put it there and show it
SPROFF n                hide it
\end{type}

\verb|SPRDEF| says where the pixels are and what shape they make:
\emph{page} is the 256-byte page the data starts at, which is how a
16-bit BASIC number reaches into 256\,MB (\verb|page*256| is the
address); \emph{w} and \emph{h} are 8, 16, 32 or 64; \emph{bpp} is 4 or
8. \verb|SPRITE| positions the sprite and turns it on, \verb|SPROFF|
hides it. Colour~0 is transparent, and there is no per-line limit --- all
128 may be on the same raster line.

Each statement re-points the chip at the attribute table and re-enables
sprites, so a \verb|GRAPHICS 0| or a shell \verb|MODE| in between costs
you nothing: the next sprite statement puts them back.

\verb|INVADER2.BAS| in \pth{/LANG/EHBASIC/EX} is the demonstration --- Space
Invaders with the arcade shapes, poked 4~bits-per-pixel into \$040000
through a bank register, four bases that erode as they are hit, and
thirty-one sprites moving. \verb|INVADERS.BAS| beside it is the same
game on the text screen, which is a fair way to see what the sprites buy
you.

\screen{shots/invaders.png}{\texttt{INVADER2.BAS}: twenty-four invaders,
four bases and a cannon, every one of them a hardware sprite. The bases
erode because their shape pages are poked, not redrawn.}

\section{The * escape}
Any shell command works from BASIC with \verb|*| in front --- \verb|*DIR|,
\verb|*CD EHBASIC|, \verb|*MON|, \verb|*DUMP a note| --- and \verb|*BYE|
leaves BASIC for the shell. Escape or Ctrl-C stops a running program (and
silences the sound); the program's variables survive for
\verb|PRINT|-style post-mortems.

\section{Editing the program in VI}\label{sec:basicvi}
\verb|*VI| and \verb|*EDIT|, with nothing after them, edit the program that
is in memory:

\begin{type}
10 PRINT "HELLO"
*VI
\end{type}

\noindent BASIC writes the program to a temp file, opens the editor on it,
and reads it back when you leave --- so what you type in the editor is what
you \verb|LIST| afterwards. Give either one a file name and nothing special
happens: \verb|*VI notes.txt| is the ordinary \verb|*| escape, and the
editor edits that file.

The trip out and back is not free. The editors load where the interpreter
lives, so the shell command underneath is \verb|SWAP|, which puts all
64\,KB and the screen aside first and restores them afterwards. And because
the program is written out and read back, \emph{variables do not survive} ---
this is an immediate-mode thing, like \verb|LOAD|. The temp file
(\pth{EDITTMP.BAS}) is left behind, which is occasionally a useful accident.

\begin{aside}
\textbf{Why the first line is blank:} \verb|SAVE| starts its output with a
newline, so the editor opens on an empty line~1 with the program below it.
Harmless --- BASIC ignores it on the way back in --- but it is why the line
count is one more than you expect.
\end{aside}

\section{The other BASIC: Microsoft, 1977}\label{sec:msbasic}
EhBASIC is the one in ROM, and the one this chapter is about. Beside it
on the disk is the other one, and it is the real thing:

\begin{type}
MSBASIC
\end{type}

\noindent Microsoft's 6502 BASIC --- not an imitation of it, but the
source Microsoft released, assembled in a \emph{pure} configuration:
the newest base in the tree, nine-digit floating point, every fix
through revision 2C, and none of the additions any particular
manufacturer made to it. It is the interpreter that became Commodore's
and Apple's before either of them touched it, and it is here because a
fantasy machine of the mid-eighties would have been sold one.

It loads at \texttt{\$7000} and comes up with its own banner. What you
should know before you type at it:

\begin{itemize}[leftmargin=1.4em]
\item \textbf{The \texttt{*} escape works}, exactly as in EhBASIC:
  \verb|*DIR|, \verb|*CD|, any shell command --- and \verb|*BYE| (or
  \verb|*QUIT|) hands the machine back to the shell with its directory
  and its screen intact. Microsoft's BASIC has no \verb|BYE| of its own;
  this one is caught a level below the interpreter, so not a byte of the
  vendored source knows it happened. A \verb|*| typed at an
  \verb|INPUT| prompt is data, not a command.
\item \textbf{It answers \texttt{MEMORY SIZE?} for you.} A pure
  Microsoft cold start asks where memory ends and, given no answer,
  finds out by walking upwards --- straight through its own image. The
  machine hands it the answer instead, and the terminal width after it,
  through BASIC's own documented input path.
\item \textbf{There is no \texttt{LOAD} or \texttt{SAVE} yet.} Both were
  things each manufacturer wrote themselves, and this build has none of
  anybody's; the machine's own are two system calls away and wiring them
  in is on the list. For now, write programs in EhBASIC, or edit a file
  with \verb|VI| and paste.
\item \textbf{It has no sound or graphics words}, and it is not going to
  grow any: adding tokens means editing the vendored interpreter, which
  is the one thing this port does not do.
\end{itemize}
